% ============ 前置部分: 记号/预备/参考文献/致谢/DRAMATIS PERSONAE ============
% 记号、预备知识与参考文献按原书连续排印, 不各自另起章节或页面.

\begingroup
\setstretch{1.0}
\clearpage
\noindent\textbf{记号.}\index{jihao@记号(notation)}

我们采用标准的(Bourbaki)记号:

$\N = \{0, 1, 2, \ldots\}$,
$\Z$ 为整数环,
$\Q$ 为有理数域,
$\R$ 为实数域,
$\CC$ 为复数域,
$\mathbb{F}_{p} = \Z/p\Z$ 为有 $p$ 个元素的域, 其中 $p$ 为素数.

对于整数 $m$ 与 $n$, $m\mid n$ 表示 $m$ 整除 $n$, 即 $n \in m\Z$. 通篇约定 $p$ 为素数, 即 $p = 2, 3, 5, \ldots$.

给定一个等价关系, $[*]$ 表示包含 $*$ 的等价类. 空集记为 $\emptyset$. 集合 $S$ 的基数记为 $|S|$ (于是当 $S$ 有限时 $|S|$ 为 $S$ 中元素的个数). 设 $I$ 与 $A$ 为集合; 以 $I$ 为指标集的 $A$ 中元素族, 记为 $(a_{i})_{i \in I}$, 是映射 $i \mapsto a_{i} : I \to A$.

$X \subset Y$ 表示 $X$ 是 $Y$ 的子集(不必是真子集);

$X \overset{\mathrm{def}}{=} Y$ 表示 $X$ 定义为 $Y$, 或按定义等于 $Y$;

$X \approx Y$ 表示 $X$ 同构于 $Y$;

$X \simeq Y$ 表示 $X$ 与 $Y$ 典范同构(或存在给定的或唯一的同构);

$\hookrightarrow$ 表示单射;

$\twoheadrightarrow$ 表示满射.

我们用哥特体(fraktur)字母表示理想, 例如
$\mathfrak{a}, \mathfrak{b}, \mathfrak{c}, \mathfrak{m}, \mathfrak{n}, \mathfrak{p}, \mathfrak{q},
 \mathfrak{A}, \mathfrak{B}, \mathfrak{C}, \mathfrak{M}, \mathfrak{N}, \mathfrak{P}, \mathfrak{Q}$
分别对应
$a, b, c, m, n, p, q, A, B, C, M, N, P, Q$.

\par\medskip\noindent\textbf{预备知识.}\index{yubeizhishi@预备知识(prerequisites)}\enspace
需要通常研究生一年级课程所覆盖的代数, 以及代数数论的一门课程, 例如我下面列出的课程讲义.

\par\medskip\noindent\textbf{参考文献.}\index{cankaowenxian@参考文献(references)}\enspace
除书末所列(以及脚注中所示的)参考文献外, 我还将引用我自己如下的课程讲义:

GT Group Theory (v3.16, 2020).

FT Fields and Galois Theory (v4.61, 2020).

ANT Algebraic Number Theory (v3.08, 2020).

\par\medskip\noindent\textbf{致谢.}\index{zhixie@致谢(acknowledgements)}\enspace

我感谢以下人士为这些讲义的早期版本提供勘误与意见: Vincenzo Acciaro; Ra\'ul Alonso Rodr\'iguez; Tom Bachmann; Oliver Braeunling; Chen, Bingxu; Kwangho Choiy; Brian Conrad; Keith Conrad; Giuseppe Canuto; Philip Dittmann; Ross Griffiths; Darij Grinberg; Florian Herzig; Chong Hui; Herv\'e Jacquet; Kiran Kedlaya; Timo Keller; Keenan Kidwell; Andrew Kirk; Michiel Kosters; Tyler Lawson; Jungin Lee; Franz Lemmermeyer; Yogesh More; Kim Nguyen; Jianing Li; Catherine O'Neil; Jack Petok; Kartik Prasanna; Nandini Ranganathan; Peter Roquette; Joshua Seaton; Corinne Sheridan; Jonah Sinick; Daniel Sparks; Jamie Tappenden; Wojtek Wawr\'ow; Yu Zhao; 以及其他人. 同时我感谢 math.stackexchange.com 与 mathoverflow.net 的贡献者, 尤其是那些在我的讲义上提问或作答的人.

我一直在读 Chevalley 关于类域论的新书; 我并非真正在做研究, 只不过是在自我修养罢了.\par\noindent Grothendieck, 1956.
\endgroup

\clearpage
\noindent\textbf{DRAMATIS PERSONAE.}\index{dramatispersonae@DRAMATIS PERSONAE(dramatis personae)}\index{renwubiao@人物表(dramatis personae)}

FERMAT (1601--1665). 提出他最后的``定理'', 并对 $m = 4$ 给出了证明. 他还提出了求方程
\begin{equation*}
X^{2} - AY^{2} = 1, \quad A \in \Z
\tag{1}
\end{equation*}
整数解的问题, 这本质上就是求 $\Z[\sqrt{A}]$ 中单位的问题\footnote{印度数学家婆什迦罗 (Bhaskara, 12 世纪) 已经知道求该方程解的一般法则。}. Brouncker 给出了求解该问题的算法, 但没有证明算法总有效.

EULER (1707--1783). 他将分析引入素数研究, 并发现了二次互反律的早期形式\footnote{Euler 发现了 [二次互反律], 而且看来有可能用散见于 Euler 的《全集》(Opera omnia) 或其遗稿中的不同片段拼出该定理的一个证明. 给出第一个完整证明的是 Gauss.$\ldots$ (Michael Berg, MR2131680).}.

LAGRANGE (1736--1813). 他证明了求解 (1) 的算法总能给出解, 并证明了每个正整数都是四个平方数之和.

LEGENDRE (1752--1833). 他引入了``Legendre 符号'' $\left(\frac{m}{p}\right)$, 并找到了二次互反律的完整形式
\[
\left(\frac{p}{q}\right)\left(\frac{q}{p}\right) = (-1)^{(p-1)(q-1)/4}, \quad p, q\ \text{为奇素数}.
\]
他证明了一个蕴含如下三元二次型局部--整体原理的结果: 对 $\Q$ 上的三元二次型 $Q(X, Y, Z)$, 它在 $\Q$ 中有非平凡零点, 当且仅当它在 $\R$ 中有, 且同余式 $Q \equiv 0 \pmod{p^{n}}$ 对所有 $p$ 与 $n$ 都有非平凡解.

GAUSS (1777--1855). 他给出了二次互反律的首批完整证明. 为寻求四次互反律, 他研究了 Gauss 整数 $\Z[\mathrm{i}]$. 他研究了 $\Z$ 上二元二次型的分类, 这与求二次域的类数问题密切相关.

DIRICHLET (1805--1859). 他引入了 $L$-级数, 并用它们证明了类数的解析公式以及算术级数中素数的密度定理. 他证明了如下的``单位定理'': 设 $\alpha$ 是首一不可约整系数多项式 $f(X)$ 的根; 若 $f(X)$ 有 $r$ 个实根与 $2s$ 个复根, 则 $\Z[\alpha]^{\times}$ 是秩为 $r + s - 1$ 的有限生成群.

KUMMER (1810--1893). 受寻求高次互反律的推动, 他对分圆域的算术作了深入研究, 并指出借助``理想数''可以恢复唯一分解. 他证明了费马大定理对正则素数成立.

HERMITE (1822--1901). 他对二次型有重要贡献, 并证明了五次多项式的根可以用椭圆函数表示. EISENSTEIN (1823--1852). 他发表了三次与四次互反律的首批完整证明.

KRONECKER (1823--1891). 他发展了 Dedekind 理想理论的一种替代方案. 他还提出了数学中最优美的想法之一, 用以生成数域的 Abel 扩张(Kronecker 青春之梦).

RIEMANN (1826--1866). 研究了 Riemann $\zeta$ 函数, 并提出 Riemann 假设.

DEDEKIND (1831--1916). 他通过给出数域中整数环的正确定义, 证明了此类环中理想可唯一分解为素理想之积, 并证明了模去主理想后它们落入有限多个类, 从而奠定了代数数论的现代基础. 他定义了数域的 $\zeta$ 函数.

WEBER (1842--1913). 他找到了``类群''允许分歧的正确推广. 在类域论与 Kronecker 青春之梦方面取得重要进展. HENSEL (1861--1941). 他在 1890 年代首次给出了 $p$-进数域的定义(作为无穷和 $\sum_{n=-k}^{\infty} a_{n}p^{n}$, $a_{n} \in \{0,1,\ldots,p-1\}$ 的集合)\footnote{赋值论由匈牙利数学家 K\"ursch\'ak 于 1912 年创立。}.

HILBERT (1862--1943). 他在 1897 年写了一本极具影响力的代数数论著作, 给出了该理论的首次系统叙述. 他的一些著名问题涉及数论, 也影响深远.

TAKAGI (1875--1960). 他证明了 Abel 类域论的基本定理, 即 Weber 与 Hilbert 所猜想者.

NOETHER (1882--1935). 她与 Artin 一道奠定了现代代数基础, 强调公理与概念性论证, 并对数域上中心单代数的分类作出贡献.

HECKE (1887--1947). 引入了 Hecke $L$-级数, 推广了 Dirichlet 的 $L$-级数与 Dedekind 的 $\zeta$ 函数.

ARTIN (1898--1962). 他发现了``Artin 互反律'', 即类域论的主定理(对 Takagi 结果的改进). 引入了 Artin $L$-级数.

HASSE (1898--1979). 他给出了局部类域论的首次证明, 证明了数域上所有二次型的 Hasse(局部--整体)原理, 并对数域上中心单代数的分类作出贡献.

BRAUER (1901--1977). 定义了 Brauer 群, 并对数域上中心单代数的分类作出贡献.

WEIL (1906--1998). 他说明了如何用 idèle 解释 Hecke 特征标. 定义了 Weil 群, 使他能够给出 Artin $L$-级数与 Hecke $L$-级数的统一推广.

CHEVALLEY (1909--1984). 类域论的主要陈述都是纯代数的, 但此前所有证明都用到了分析; Chevalley 给出了纯代数证明. 借助他引入的 idèle, 他得以对无穷 Abel 扩张给出类域论的自然表述.

IWASAWA (1917--1998). 他给代数数论引入了一种重要的新方法, 其灵感来自有限域上曲线的理论.

TATE (1925--2019). 他在群上同调中证明了新结果, 使他能对类域论给出优雅的重新表述. 他与 Lubin 一起找到了生成局部域 Abel 扩张的一种显式方法.

LANGLANDS (1936--). Langlands 纲领\footnote{不要与它的几何类比相混淆, 后者有时称为几何 Langlands 纲领, 它似乎缺乏算术与解析意义。} 是一系列宏大的猜想,其中除其他内容外, 包含了一个非 Abel 的类域论.
